% Generated by scripts/tikz_reviews.py using the compiled book's labels.
\pdfbookmark[0]{Chapter 3}{comparison-chapter-3}
\ReviewFigure{fig-entropy-extrema}
{Entropy level sets and probability constraints}
{figures/1-shannon/entropy-extrema.pdf}
{figures/tikz/coding-entropy-extrema.pdf}
{The entropy lemma states that point masses minimize entropy and the uniform distribution maximizes it.}
{The new left panel shows actual level sets of -p1 log2(p1)-p2 log2(p2) in the positive quadrant, with 0 log2(0)=0. This formula extends beyond probability vectors; the red segment p1+p2=1 is the probability constraint. Its constrained maximum is at (1/2,1/2), where H=1 and the contour is tangent to the segment. The unconstrained maximum of the extension is instead at (1/e,1/e), explaining why the contours are centered off the constraint. The middle panel restricts entropy to the segment. The right panel retains the three-symbol simplex, where the uniform entropy is log2(3).}
{coding-entropy-extrema}
{Complete figure}
{3.4}
\ReviewFigure{fig-kraft-trees}
{Kraft inequality: disjoint descendant blocks}
{figures/1-shannon/kraft-ineq-1.png}
{figures/tikz/coding-kraft-necessity.pdf}
{The necessity proof completes a prefix tree to depth m and counts the descendants of each codeword.}
{The code is (0,10,110,111). Its codewords reserve 4, 2, 1, and 1 leaves in the full depth-three tree. Tinted subtrees connect the highlighted codeword roots to disjoint descendant blocks. Edge bits show how a path spells a codeword; the formula explains the count 2\textasciicircum{}(m-l).}
{coding-kraft-necessity}
{Panel 1 of 2}
{3.5}
\ReviewFigure{fig-kraft-trees}
{Kraft inequality: packing prescribed lengths}
{figures/1-shannon/kraft-ineq-2.png}
{figures/tikz/coding-kraft-sufficiency.pdf}
{The converse proof packs dyadically aligned blocks in nonincreasing order of size.}
{For prescribed lengths (1,2,3), the blocks contain 4, 2, and 1 leaves. The left-to-right packing arrow and tinted subtrees expose the construction: choose the largest block first, then continue at the next dyadic boundary. This gives codewords (0,10,110), leaves 111 unused, and demonstrates 7/8 <= 1 without assuming equality.}
{coding-kraft-sufficiency}
{Panel 2 of 2}
{3.5}
